\section{The inner radius and the local-time profile}\label{sec:contact}
We compare the range with its largest centered ball and then determine
the radius from the total elapsed time.

Fix $p>0$ and intersect the events in Proposition~\ref{prop:coarse},
Lemmas~\ref{lem:local} and~\ref{lem:radial}, and
\eqref{eq:localmart} and~\eqref{eq:vector}. The complement has probability
at most $Cn^{-p}$. All arguments below are deterministic on this event,
apart from the additional concentration estimate for the quadratic
martingale, whose exceptional probability has the same bound.
Choose a deterministic $K>a+1$ large enough that the coarse radius
bound implies $D_n\subset B(0,(K-1)N)$. On $B(0,KN)$,
\eqref{eq:approx} gives
\begin{equation}\label{eq:global}
|\widetilde\ell_n(y)-U_{D_n}(y)|\leq\delta_0 ,\qquad
\delta_0=\begin{cases}C\sqrt N L,&d=2,\\
C(\sqrt{NL}+L),&d\geq3
\end{cases}\, .
\end{equation}


Let $D=D_n$ and let $W=NQ^{1/d}$. Write the centered inradius and excess volume as
\begin{equation}\label{eq:bm}
 b=\inf\{|y|:y\notin D\},\qquad E=D\setminus B(0,b),\qquad m=|E|\, .
\end{equation}
The ball $B(0,b)$ is contained in $D$, up to immaterial cell boundaries.
Choose $y_0$ of norm $b$ that is a limit of points outside $D$.
There is an unoccupied lattice cell $C_z$ whose closure contains $y_0$.
Thus
\begin{equation}\label{eq:contactcell}
 \ell_n(z)=0,\quad |z-y_0|\leq\sqrt d/2,\quad |z|\geq b\, .
\end{equation}
This follows by taking a convergent subsequence from the finitely
many cells adjacent to $y_0$; no regularity of $\partial D$ is assumed.

\subsection{The potential at an unvisited site}
The excess volume gives a positive contribution to the potential at $y_0$.
We claim that
\begin{equation}\label{eq:contactbound}
c\frac{m}{N^{d-1}}\leq U_{D_n}(y_0)
\leq C(\sqrt{B_zL}+L)\, .
\end{equation}
In addition, for every $y\in\R^d$,
\begin{equation}\label{eq:packing}
|U_E(y)|\leq Cm^{1/d}\, .
\end{equation}

If $s=|v|\geq b=|y_0|$, then
\[
 \frac{u_v\cdot(v-y_0)}{|v-y_0|^d}
 =\frac{s^2-v\cdot y_0}{s|v-y_0|^d}
 \geq\frac1{2s}|v-y_0|^{2-d}
 \geq 2^{1-d}s^{1-d}\, .
\]
The first inequality uses
$2(s^2-v\cdot y_0)=|v-y_0|^2+s^2-b^2$; the second uses
$|v-y_0|\leq2s$ and $d\geq2$.
The singular point is immaterial in the integral. Since
$U_{B(0,b)}(y_0)=0$ and $s\leq CN$ on $D$, integration proves
the lower bound in \eqref{eq:contactbound}. The upper bound follows
from \eqref{eq:pointwise}, \eqref{eq:localmart}, \eqref{eq:cellmodulus}, and
$\ell_n(z)=0$; the target is admissible because $|z|=O(N)<3n$. Finally $|u_v|\leq1$, and integrating
$|v-y|^{1-d}$ over a set of volume $m$ is maximized by a centered ball.
Its integral is $O(m^{1/d})$, proving \eqref{eq:packing}.



\subsection{The contact variance in higher dimensions}
Suppose that $d\geq3$. A summable convolution bounds the local variance
near the unvisited site.
Put $V=m^{1/d}$ and $s_b(y)=(b-|y|)_+$. Equations
\eqref{eq:pointwise}, \eqref{eq:localmart}, \eqref{eq:ballpotential}, and
\eqref{eq:packing} imply, at occupied lattice sites,
\[
 \ell_n(y)\leq c_0s_b(y)+C(V+L)+C\sqrt{B_yL}\, .
\]
Let
\[
 k(x)=(1+|x|)^{2-2d}\ind_{\{|x|\leq3n\}},\quad
 k_0=\sum_xk(x),\quad J=\sum_x|x|k(x)\, .
\]
For sufficiently large $n$, \eqref{eq:coarse} ensures that
$B_y=(k*\ell_n)(y)$ at every occupied site and at $z$.
Moreover
\begin{equation}\label{eq:kernelmoments}
 k_0\leq C,\qquad
 J\leq\begin{cases}CL,&d=3,\\C,&d\geq4\end{cases}\, .
\end{equation}
Choose a deterministic $\lambda>0$ so small that
$\lambda k_0<1/2$ for every $n$. Young's inequality gives
\begin{equation}\label{eq:resolvent}
 \ell_n\leq c_0s_b+C(V+L)+\lambda k*\ell_n
 \quad\hbox{on all of }\Z^d\, .
\end{equation}
Outside the occupied set the inequality holds trivially because its
left side is zero and its right side is nonnegative.

Since $s_b$ is $1$-Lipschitz,
\[
 (k^{*j}*s_b)(y)\leq k_0^j s_b(y)+j k_0^{j-1}J\qquad(j\geq1)\, .
\]
Iterate \eqref{eq:resolvent} and sum the two geometric series. The remaining
term tends to zero since the operator norm on $\ell^\infty$ is
$\lambda k_0<1/2$ and $\ell_n$ is bounded.
Using \eqref{eq:kernelmoments} yields
\begin{equation}\label{eq:envelopehigh}
 \ell_n(y)\leq C(s_b(y)+V+L)\quad(y\in\Z^d)\, .
\end{equation}
At the unoccupied center $z$, $s_b(z)=0$, so convolution once more gives
\begin{equation}\label{eq:holebracket}
 B_z\leq C(V+L)\, .
\end{equation}
Indeed $s_b(x)\leq|x-z|$ and hence $(k*s_b)(z)\leq J$.
The contact inequality now becomes
\[
 m\leq CN^{d-1}\big(m^{1/(2d)}\sqrt L+L\big)\, .
\]
If the second term is at least half the left side, then
$m\leq CN^{d-1}L$. Otherwise raise the inequality for
$m^{1-1/(2d)}$ to the power $2d/(2d-1)$. In both cases,
\begin{equation}\label{eq:masshigh}
 m\leq CN^{2d(d-1)/(2d-1)}L^{d/(2d-1)}
       =CN^d Q\, .
\end{equation}
The bound $N^{d-1}L$ is absorbed because $N/L\to\infty$.

\subsection{The planar contact variance}
Suppose that $d=2$. A first bound on the excess volume removes the
logarithmic loss in the variance at the unvisited site.
The coarse global approximation \eqref{eq:global} and the contact
inequality first give
\[
 m\leq CN^{3/2}L\, .
\]
By \eqref{eq:ballpotential} and \eqref{eq:packing}, for every occupied $x$,
\[
 \ell_n(x)\leq c_0(b-|x|)_+
       +CN^{3/4}L^{1/2}+C\sqrt N L\, .
\]
At the unoccupied center $z$, $(b-|x|)_+\leq|x-z|$.
All occupied sites lie within distance $CN$ of $z$, so
\[
 \sum_{|x-z|\leq CN}\frac{(b-|x|)_+}{(1+|x-z|)^2}
 \leq CN,\qquad
 \sum_{|x-z|\leq CN}(1+|x-z|)^{-2}\leq CL\, .
\]
Therefore
\[
 B_z\leq C\big(N+N^{3/4}L^{3/2}+\sqrt N L^2\big)
       \leq CN\, .
\]
for all sufficiently large $n$. Applying
\eqref{eq:contactbound} again proves
\begin{equation}\label{eq:massplanar}
 m\leq CN^{3/2}\sqrt L=CN^2Q\, .
\end{equation}
Combining this with \eqref{eq:global} and \eqref{eq:packing} gives
\begin{equation}\label{eq:envelopeplanar}
 \ell_n(y)\leq c_0(b-|y|)_++CW,
 \qquad W=N^{3/4}L^{1/4}\, .
\end{equation}
Here $\sqrt N L=o(W)$.
In every dimension, \eqref{eq:envelopehigh} or \eqref{eq:envelopeplanar}
therefore implies
\begin{equation}\label{eq:shellW}
 M_{\rm sh}(r)\leq CW\quad\hbox{whenever }r\geq b+3\, .
\end{equation}

\subsection{Determining the radius}
The squared displacement determines the radius through the elapsed time.
Since every step has squared length one and conditional mean
$-\eps I_ju_{X_j}$, the process
\[
\mathcal Q_t=\sum_{j<t}2X_j\cdot(X_{j+1}-X_j+\eps I_ju_{X_j})\, ,
\]
is a martingale, and
\begin{equation}\label{eq:quadratic}
 |X_n|^2=n-2\eps\sum_{x\in A_n}|x|+\mathcal Q_n\, .
\end{equation}
Stop this martingale upon first exiting $B(0,KN)$. Its increments are bounded by $CN$ and its
bracket by $CnN^2$. Freedman's inequality, intersected with the
coarse event, yields
\begin{equation}\label{eq:quadraticerror}
 |\mathcal Q_n|\leq C\big(N\sqrt{nL}+NL\big)\, .
\end{equation}
with another failure probability at most $Cn^{-p}$.
On the coarse event the stopped and original martingales coincide.

The cell replacement satisfies
\[
 \left|\sum_{x\in A_n}|x|-\int_D|v|\dd v\right|
 \leq (\sqrt d/2)|A_n|\leq CN^d\, .
\]
Also
\[
 \int_D|v|\dd v
 =\frac{d\omega_d}{d+1}b^{d+1}+\int_E|v|\dd v,
 \qquad 0\leq\int_E|v|\dd v\leq CNm\leq CN^{d+1}Q\, .
\]
Divide \eqref{eq:quadratic} by $n=N^{d+1}$. Equations
\eqref{eq:quadraticerror}, \eqref{eq:masshigh}, and \eqref{eq:massplanar} give
\[
 \left|1-\frac{c_0\omega_d}{d+1}(b/N)^{d+1}\right|
 \leq C\left(Q+N^{-1}+N^{-(d-1)/2}\sqrt L
                         +N^{-d}L+N^{1-d}\right)
 \leq CQ\, .
\]
For $d=2$, the largest additional term is $N^{-1/2}\sqrt L=Q$.
For $d\geq3$, it is $O(N^{-1}\sqrt L)=O(Q)$.
Since the positive root of
$c_0\omega_d t^{d+1}/(d+1)=1$ is $a$, this proves
\begin{equation}\label{eq:inradius}
 |b/N-a|\leq CQ\, .
\end{equation}
In particular $b\asymp N$. It follows immediately from
$D=B(0,b)\cup E$ up to null sets that \eqref{eq:volume} holds.
The inner inclusion in \eqref{eq:sandwich} follows by increasing its
constant to avoid cell-boundary equality.

Finally, on the approximation region, \eqref{eq:global},
\eqref{eq:ballpotential}, \eqref{eq:packing}, and \eqref{eq:inradius} give
\[
 \left|\widetilde\ell_n(y)-c_0(aN-|y|)_+\right|
 \leq C(\delta_0+NQ+NQ^{1/d})\leq CW\, .
\]
This proves \eqref{eq:profile-rate}; outside a fixed ball containing the range
both the local time and the comparison cone vanish.

\section{The outer fluctuation bound}\label{sec:outer}
We bound the weighted exterior volume and use the crossing estimate
to exclude excursions beyond the claimed radius.
Start at an integer radius just above $b+2b_d+6$. Since $b\asymp N$,
\[
 F(b)\leq CmN^{1-d}\leq CNQ\, .
\]
All subsequent radii remain in a fixed annulus of scale $N$.
Let $\Lambda=CN^{2-d}L$, choosing $C$ sufficiently large.
If a current dyadic upper bound $A$ satisfies $A\geq\Lambda$, then the
error in \eqref{eq:radial}, when $F(r-b_d)\leq A$, is at most $A/4$.
If $F(r+b_d)>A/2$, then \eqref{eq:radial} and \eqref{eq:shellW} therefore force
\[
 F(r-b_d)-F(r+b_d)\geq \frac{cA}{W+1}\, .
\]
On the grid of spacing $2b_d$, a halving takes at most $C(W+1)$
steps. There are at most $CL$ levels between $CNQ$ and
$\Lambda$. If the initial bound is already below $\Lambda$, then no
halving is needed. Thus, for a deterministic constant $K_0$,
\begin{equation}\label{eq:tailend}
 F(b+K_0WL)\leq\Lambda\, .
\end{equation}
Rounding errors are absorbed because $W\to\infty$.
Since $WL=o(N)$, all required radii fit in that fixed annulus.

Set $h=K_1WL$, with $K_1>2K_0$, and suppose that the path
reaches radius $r=b+3h$ by time $n$. Let
\[
 \tau=\min\{t\leq n:|X_t|\geq r\},\quad
 v=X_\tau/|X_\tau|,\quad \beta=b+2h,
 \quad s=1+\max\{j<\tau:v\cdot X_j\leq\beta\}\, .
\]
The set defining $s$ contains zero. All departures from $s$ through
$\tau-1$ have projection greater than $\beta$ and radius less than
$r$, while their net projected displacement is at least $h-1$.
The ratio $\beta/r$ is bounded below by a fixed positive constant.
The distinct departure cells in this interval lie in $D_n$ even
if $\tau=n$. By \eqref{eq:tailend}, their number $m_{\rm cap}$ satisfies
\[
 m_{\rm cap}\leq CN^{d-1}F(\beta-\sqrt d/2)
              \leq CNL\, .
\]
Consequently \eqref{eq:crossing} would imply
\begin{equation}\label{eq:outer-contradiction}
 (h-1)^2\leq C
 \begin{cases}
 N^{4/3}L^{8/3}+NL^4,&d=2,\\
 N^{2d/(2d-1)}L^{4d/(2d-1)}+NL^3,&d\geq3
 \end{cases}\, .
\end{equation}
The right side is $o(h^2)$: for $d=2$,
$h^2\asymp N^{3/2}L^{5/2}$; for $d\geq3$,
$h^2\asymp N^{(4d-4)/(2d-1)}L^{4d/(2d-1)}$ and
$4d-4>2d$. This contradiction proves the outer inclusion in
\eqref{eq:sandwich}, using \eqref{eq:inradius} and $NQ=O(WL)$.

The lattice mesh contributes at most $\sqrt d/N$ to Hausdorff distance,
which is absorbed in \eqref{eq:hausdorff}.



\begin{proof}[Proof of Theorems~\ref{thm:shape} and~\ref{thm:fluctuations}]
The preceding estimates prove the inequalities in
Theorem~\ref{thm:fluctuations} at each deterministic $n$, with total
failure probability at most $Cn^{-p}$. Fix $p>1$. Borel--Cantelli
implies that these inequalities hold for every sufficiently large
integer $n$, almost surely. Since $Q\to0$ and $Q^{1/d}L\to0$,
the ball inclusions and profile convergence in Theorem~\ref{thm:shape}
follow. The volume estimate gives $|V_n|/N^d\to\omega_da^d$ because
$|V_n\setminus A_n|\leq1$. For each fixed lattice site $x$,
the uniform profile estimate and $|x|/N\to0$ give
$\ell_n(x)/N\to2d\eps a>0$. The lattice is countable, so the
walk visits every vertex infinitely often on the same probability-one event.
\end{proof}
